\section{K-038: the odd-root component and its local moving coefficient}

The finite quotient in K-037 is in fact the normalization of an
actual three-dimensional odd-root fixed component. The additional
input is an exact unitary covariant which controls the full
stabilizer of the displayed critical tensor. We then evaluate its
relative moving Euler factor. This is local component data; it
does not evaluate the specific dyadic exceptional class of K-031.

\subsection*{The full stabilizer at the displayed point}

Retain the critical dimension-eighteen tensor $T$, the dyadic
element $h$, the cubic element $a$, and $g=ha$ from K-037.
Write $\widetilde H=\operatorname{Stab}_{GL_{18}^3}(T)$,
$S=\{(\alpha I,\beta I,\gamma I):\alpha\beta\gamma=1\}$,
and $H=\widetilde H/S$.

We first determine $J=C_{\widetilde H}(h)$.
Criticality gives adjoint closure by the polar argument in K-035.
All old frames are nonzero, so $J$ preserves the boundary line
and its orthogonal two-plane in each selected three-block.
The projection of $T$ onto
$A_{\rm new}\otimes B_0\otimes C_{\rm new}$ is the nonzero
rank-one tensor $e_2\otimes e_0\otimes e_2/\sqrt3$.
Thus $J$ preserves the second new line in $A$ and $C$.
The analogous projection with $C_0$ preserves that line in
$B$ as well. Adjoint closure preserves their complementary
new lines. Therefore $J$ is entirely diagonal.

The boundary coefficient $\sqrt{\rho-1/3}$ is nonzero,
so the old coordinates satisfy every old support equation.
Their diagonal stabilizer is the already proved
$S_{16}\times\mu_D$. Divide each new eigenvalue by its old
boundary eigenvalue and denote the six ratios by
$a_1,a_2,b_1,b_2,c_1,c_2$. The seven nonboundary core equations
are
\[
 b_1c_2=b_2c_1=a_1b_1c_2=a_1b_2c_1
          =a_2c_2=a_2b_1c_1=a_2b_2=1.
\]
They are equivalent to
\[
 a_1=1,\qquad a_2=b_1=c_1=\theta,\qquad
 b_2=c_2=\theta^{-1},\qquad\theta^3=1.                    \tag{K-038a}
\]
Their determinant-product ratio is $\theta$.
Consequently $J/S=\langle h,a\rangle\simeq\mu_{3D}$,
and its determinant character is faithful.

The effective centralizer $C_H(\bar h)$ equals $J/S$:
a scalar shift preserving an $h$ spectrum must fix its
unique triple eigenvalue and is one. Also $H$ is finite.
Indeed the stable-maximal-torus proof of K-035b applies
in any dimension and with any eigenvalue multiplicities:
the first and third $h$ spectra are antipodal-free, and
the representation on those two natural factors is
faithful by conciseness. If $H$ had positive dimension,
that proof would give a positive-dimensional $J/S$,
contrary to (K-038a).

We now show that every full stabilizer centralizes $h$.
For a tensor $R$ whose first factor is the factor being
tested, define the matrix covariant
\[
 X(R)_{u,u'}=
 \sum R_{u b_1c_1}R_{u_2b_2c_2}R_{u_3b_3c_3}
 \overline{R_{u'b_2c_3}R_{u_2b_3c_1}R_{u_3b_1c_2}}.
                                                               \tag{K-038b}
\]
All repeated indices are summed. Direct contraction of
unitary matrices shows
\[
 X((k_1,k_2,k_3)R)=k_1X(R)k_1^*.
\]
Apply this after placing factor $C$ first. The exact
calculation on $T$ gives a diagonal matrix with eighteen
distinct rational diagonal entries. There are only
$224$ nonzero summands in (K-038b). The full list of
entries and all contraction terms are retained in
\nolinkurl{rank_three_full_stabilizer_checks.json}.
Distinctness is an exact rational comparison, not a
numerical spectral test.

Since $H$ is finite and $T$ critical, every effective
stabilizer class has a unitary representative: in its
polar decomposition the positive factor lies in the
identity component $S$. For such a representative $k$,
(K-038b) implies that $k_C$ commutes with the displayed
simple-spectrum diagonal matrix. Hence $k_C$ is diagonal
in the same coordinates and commutes with $h_C$.

Here is the needed injectivity of $H\longrightarrow PGL(C)$.
An element in its kernel has scalar third factor.
Multiplying by an element of $S$ makes that factor the
identity. Write the slices of $T$ as matrices $M_j$ in
the other two factors. The pointwise-third-factor
stabilizer is the unit group of the finite-dimensional
endomorphism algebra
\[
                 \mathcal A=\{(u,v):uM_j=M_jv
                                      \text{ for every }j\},
                                                               \tag{K-038c}
\]
under $(u,v)\mapsto(u,v^{-t},I)$.
Its unit group is a nonempty Zariski-open subset of the
vector space $\mathcal A$, and is therefore connected.
Its image is exactly the indicated kernel in the finite
group $H$. That kernel is consequently trivial.
No invertible-slice assumption is used.

The commutator of $k$ and $h$ projects trivially to
$PGL(C)$ and is therefore trivial in $H$. Thus
$H=C_H(\bar h)=J/S$. We have proved
\[
        \operatorname{Stab}_{GL_{18}^3}(T)
                  =S\times\mu_{3D},
          \qquad\chi\text{ faithful on }\mu_{3D}.        \tag{K-038d}
\]
In particular the effective special-linear stabilizer is
trivial, and the tensor quotient is smooth at its image.

\subsection*{The actual odd-root component and its normalization}

At this point the $g$-fixed quotient tangent space has
dimension $56-53=3$, by K-037j and exactness of finite-order
fixed vectors in the tangent sequence. Surjectivity of
the fixed differential, followed by the argument of
K-029, identifies the closed irreducible image $Z$ of
$U_g$ with a whole irreducible component of
$X_{18}^{\mu_{3D}}$. It is smooth at the displayed point.

The finite map $\operatorname{Spec}B\to Z$ in K-037 is
birational. On the nonempty effectively free open set,
the full effective stabilizer has injective character.
Two closed $g$-fixed representatives of the same quotient
point are special-linearly related; their two conjugate
lifts of the same grading character therefore agree
modulo $S$. There are no nontrivial scalar spectral
shifts of $g$ in any factor. To prove this, note that
$g^3=h^3$ has a unique triple eigenvalue and fifteen
simple eigenvalues in each factor. A scalar shift
$\beta$ preserving the $g$ spectrum must satisfy
$\beta^3=1$ by the unique triple eigenvalue of $g^3$.
If $\beta$ were a nontrivial cube root, its rotation
would have to preserve each one-dimensional $g^3$
eigenspace, which is impossible. Thus $\beta=1$.
The relating element centralizes the actual $g$ and
belongs to $C_g$.

The finite map consequently has one generic geometric
point in its fibre and degree one in characteristic zero.
The invariant ring $B$ is normal, so it is the
normalization of $Z$. Its density is therefore also
the top density of $Z$:
\[
                            e_Z=9063377350               \tag{K-038e}
\]
in the grading $z_3=t^{3D}$. The conductor, if present,
has smaller dimension and cannot affect that top term.

\subsection*{Moving characters and the precise local parameter}

Fix a primitive $3D$th root $\zeta$, write
$\omega=\zeta^D$, and realize $g$ with the exponent
lists of K-037's certificate. Let
$\xi=\chi(g)$ be the corresponding primitive grading
root. We use the point-pullback convention of P-F23
and K-014: the eigenvalues in the moving inverse Euler
factor are those of the displayed fixing element on
the tangent space.

The $h$-fixed tangent space has dimension thirteen.
Within it the nontrivial cubic tangent eigenspaces
have multiplicities
\[
                     \omega:6,\qquad\omega^2:4.          \tag{K-038f}
\]
Indeed the corresponding multiplicities in $V$ are
$14,12$, and in the special-linear Lie algebra they
are $8,8$. Hence the relative inverse Euler factor
inside the thirteen-dimensional component $Y$ is
\[
 E_{\rm rel}
   =(1-\omega)^{-6}(1-\omega^2)^{-4}
                    =-\frac{\omega^2}{3^5}.             \tag{K-038g}
\]
All other ambient moving tangent eigenvalues have
nontrivial dyadic part. Their factors $1-\lambda$
are units at three. Thus the full moving factor is
\[
                 E_g=E_{\rm rel}\,U_g,\qquad
                 U_g\text{ a three-adic unit},           \tag{K-038h}
\]
where the exact finite product for $U_g$ is specified
by the full moving-tangent multiplicity list in the
certificate. In particular $v_3(E_g)=-5$.

The local leading contribution of this actual
three-dimensional component at $t=\xi$ is
\[
                     \frac{e_ZE_g}{(1-t^{3D})^3}
               +\text{terms of lower pole order}.       \tag{K-038i}
\]
Its leading coefficient in the stated parameter
$v=1-t^{3D}$ has valuation $-5$ at three.
Inside $Y$ itself, the corresponding coefficient is
$e_ZE_{\rm rel}$ in that same parameter.

These normalizations matter. If the local parameter
is instead $w=1-t/\xi$, then
$v=3Dw+O(w^2)$, so the coefficient of $w^{-3}$ is
$e_ZE_g/(3D)^3$ and has valuation $-8$.
For the grading of $Y$, put $z=t^D$. At this point
$z=\xi^D=\omega^2$, since $D=2^{17}\equiv2\pmod3$.
The relative local parameter is $1-z^3$, and
$E_{\rm rel}=-z/3^5$ at that root.
None of these statements is the expansion at the
dyadic root used in K-037s, where $1-t^D$ vanishes
and the supported-module coefficient had valuation
$-3$. They are different local expansions.

The local coefficient (K-038i) is nonzero. It does
not prove that the sum over all odd-root fixed
components is nonzero, nor does it evaluate the
specific dyadic collision class of K-031.
The full thirteen-dimensional density of $Y$ and
its contribution at the primitive dyadic root remain
unknown.

\paragraph{Exact check.}
The verifier
\nolinkurl{verify_rank_three_full_stabilizer.py}
reconstructs every contraction in (K-038b), checks
the eighteen distinct rational eigenvalues, verifies
the six-variable diagonal equations and their index
three, checks all scalar spectral periods, and
computes the full ambient tangent character and
the multiplicities (K-038f).
The first execution completed all mathematical
assertions but encountered a JSON serialization
error on SymPy integers; conversion to ordinary
integers was corrected and the complete replay
passed. No mathematical assertion was weakened.

